\section{Discussion, Limitations, and Conclusion}
\label{sec:limitations}

\paragraph{Correlational evidence is descriptive.} Section~\ref{sec:rq1}'s correlations use four $\lambda$ points per dataset; the \emph{sign reversal} between datasets is the finding, not $n=4$ significance.

\paragraph{FashionIQ's small Recall is not noise.} All FashionIQ R@1 values are below $1\%$, but the vanilla-vs-regularized gap holds at R@5/R@10 too, 3-seed mean $\pm$ std (vanilla $2.94\pm0.02\%$/$4.76\pm0.05\%$ vs.\ best regularized, weak, $1.16\pm0.14\%$/$1.82\pm0.15\%$) -- not a $k{=}1$ artifact. The ranking \emph{among} regularized variants is not seed-stable; we claim none.

\paragraph{Seed 2023's original eval was an artifact in R@1, not only R@5/R@10.} Seed 2023's R@5/R@10 originally showed near-flat growth over R@1 (vanilla: $14.08\to18.08\to18.11\%$) vs.\ seeds 7/13/21/42's normal $\sim$2--2.5$\times$ growth -- a truncated ranked list, not a genuine effect. Re-evaluating that point (identical checkpoint and config) also corrected R@1: vanilla moved $14.08\%\to20.86\%$, now consistent with its other four seeds ($20.70$--$20.81\%$) rather than a $\sim$6-point outlier that drove this paper's original $+50.9\%\to+22.0\%\to$``not significant'' narrowing (Section~\ref{sec:rq2rq3}). We ruled out the two likely causes: the checkpoint is unchanged (timestamp-confirmed), and the reused trie passed a hash check against current codes -- not the documented stale-cache bug. The remaining explanation is genuine run-to-run non-determinism in beam-search decoding; disclosed as open, not solved. Corrected values are used throughout (Tables~\ref{tab:table2}--\ref{tab:localization}, Figure~\ref{fig:lambda-recall}).

\paragraph{Remaining open points, and scope.} The I$\to$T collapse is localized to the tokenizer (Section~\ref{sec:rq2rq3}) and partially explained by text embeddings' greater anisotropy under moderate regularization, but not why it persists undiminished at strong $\lambda$ once code entropy partially recovers -- a gap we disclose rather than paper over. The dense-retrieval ceiling of Section~\ref{sec:setup} bounds all variants here, sitting 3--8$\times$ below it and the released model. Our datasets are structurally contrasting exemplars, not a sweep; the optimal $\lambda$ found need not transfer to a further regime. Adapting Wang et al.'s decoupled-RQ design~\cite{wang2026decoupledrq} to our cross-modal setting is out of scope. VisualNews was tested only at $\lambda\in\{0,3\}$; a favorable intermediate strength cannot be excluded.

\paragraph{Direction-aware $\lambda$'s decode crash: nondeterministic, not structural, and now patched.} Constrained-decoding evaluation for img3txt0/img3txt0p3 originally failed 12/12 attempts across three retrain-and-eval rounds, all traced to the same root cause: HuggingFace's constrained beam search reaching a decode state with zero valid trie continuations (10/12 as the synchronous \texttt{prefix\_allowed\_tokens\_fn returned an empty list}; the other 2 as a deferred CUDA assertion from the identical call path). These are the two most level-1-concentrated checkpoints trained (603/391 unique codes vs.\ vanilla's 1,249), so severe codebook collapse plausibly makes this dead-end more likely. We patched the decoder to fall back to the trie's root-level tokens for exactly this case, then re-ran evaluation on the existing checkpoints for two additional seeds; all four runs completed via unmodified beam search, with the fallback never triggering -- itself suspicious for a claimed-deterministic bug on byte-identical checkpoints. We tested directly: re-running the \emph{original, unpatched} code twice on the same four checkpoints (8 reruns) produced zero crashes, despite 12/12 original failures on these exact weights. Checkpoint and trie were confirmed byte-identical throughout, so the only reconciliation is genuine run-to-run nondeterminism: constrained beam search runs in fp16 with no deterministic CUDA/cuDNN kernels, so on a severely collapsed codebook with many near-tied beam candidates, floating-point differences can change which beam hits the empty-continuation state -- recurring across three independent \emph{training} rounds (each producing fresh near-ties) yet not across repeated \emph{evaluation} of one fixed checkpoint. The fallback still converts an unpredictable crash into a guaranteed-terminating decode, and Table~\ref{tab:direction-aware}'s numbers are unaffected either way, since it never fired in the runs that produced them. The qualitative finding holds: both variants underperform vanilla on T$\to$I and stay collapsed on I$\to$T, tight across seeds ($7.87\pm0.12\%$/$5.56\pm0.13\%$ T$\to$I; $0.12\pm0.04\%$/$0.13\pm0.05\%$ I$\to$T), matching Section~\ref{sec:direction-aware}'s tokenizer-level probe.

\paragraph{Conclusion.} Balance regularization is neither useless nor universally good: it significantly improves MSCOCO T$\to$I ($+12.3\%$, $p<10^{-6}$, verified across both decoder seeds and independently retrained tokenizers) while destroying I$\to$T on MSCOCO and VisualNews, substantially \emph{hurts} VisualNews's T$\to$I, and hurts FashionIQ outright. Decoupling per modality does not rescue either direction, regardless of which side carries the weight. Do not expect MSCOCO's gain to transfer.

Code and per-seed results will be released as an anonymized repository at submission time.
